\section{Требования к программе}

\subsection{Функциональные требования}

\headingtextsep

Программная система должна обеспечивать следующие функциональные возможности.

\subsubsection{Рабочие пространства, проекты и участники}

Система должна поддерживать создание нескольких рабочих пространств.

Внутри каждого рабочего пространства система должна поддерживать создание и ведение проектов.

Система должна обеспечивать приглашение участников, назначение ролей, отключение участника и передачу прав владения рабочим пространством.

Система должна обеспечивать разграничение доступа между рабочим пространством, проектом и административными разделами.

\subsubsection{Задачи, группы задач и зависимости}

Система должна обеспечивать создание, редактирование, архивирование и поиск задач.

Система должна поддерживать статусы, приоритеты, сроки, исполнителя, иерархию подзадач и порядок отображения задач.

Система должна обеспечивать группировку задач в логические блоки работ, в том числе инициативы, эпики и иные крупные группы.

Система должна обеспечивать хранение направленных зависимостей между задачами и определение блокирующих связей на их основе.

Система должна обеспечивать представления списка, доски, иерархии и иных видов отображения на основе единой модели задач.

Система должна позволять назначать задачу как человеку-участнику, так и программному агенту.

\subsubsection{Документация, заметки и файловые ресурсы}

Система должна обеспечивать ведение проектной базы знаний в виде страниц с содержимым в формате Markdown.

Система должна поддерживать древовидную структуру документов, внутренние ссылки, версии содержимого и ссылки на связанные файловые ресурсы.

Система должна обеспечивать создание и хранение заметок, связанных с задачами, документами, проектом или агентной сессией.

Система должна обеспечивать загрузку и хранение файловых ресурсов отдельно от документов с возможностью последующего связывания их с сущностями проекта.

\subsubsection{Агентный доступ и машинное взаимодействие}

Система должна поддерживать регистрацию программных агентов в пределах рабочего пространства.

Система должна обеспечивать выпуск, отзыв, ротацию и ограничение областей действия учётных данных агентов.

Система должна фиксировать сессии работы агентов, их действия и затронутые задачи.

Система должна предоставлять машинный доступ к части функций через программный интерфейс и MCP-совместимый интерфейс.

\subsubsection{Поиск, интеграции и аудит}

Система должна обеспечивать полнотекстовый и семантический поиск по задачам, документам, заметкам, метаданным файловых ресурсов и результатам агентной активности.

Система должна обеспечивать хранение и просмотр журнала значимых действий пользователей, агентов и интеграций.

Система должна обеспечивать интеграцию с внешними сервисами, в том числе с GitHub, для связывания задач с issue, pull request, commit и branch.

Система должна поддерживать хранение унифицированных ссылок между внутренними сущностями системы и внешними ресурсами.

\subsubsection{Режимы функционирования}

Основной режим:

\begin{itemize}
    \item выполнение полного набора пользовательских, административных и агентных операций;
    \item хранение проектных данных, документов, файловых ресурсов и событий аудита;
    \item работа интеграций, поиска и механизмов машинного доступа.
\end{itemize}

Локальный режим разработки:

\begin{itemize}
    \item запуск системы с упрощенной схемой аутентификации для отладки;
    \item использование локального файлового хранилища;
    \item возможность запуска без внешних облачных учётных данных;
    \item наличие тестовых участников и тестового агента, доступных только в локальной среде разработки.
\end{itemize}

\textheadingsep
\subsection{Требования к интерфейсу}

\headingtextsep

\subsubsection{Пользовательский веб-интерфейс}

\paragraph{Адаптивность}

Веб-интерфейс должен корректно работать на экранах настольных и мобильных устройств в пределах поддерживаемых разрешений.

\paragraph{Локализация}

Система должна поддерживать русский язык интерфейса.

\paragraph{Общие принципы оформления}

Интерфейс должен быть строгим, функциональным и ориентированным на работу с насыщенными информацией экранами.

Основная рабочая среда должна быть оптимизирована для настольных устройств; мобильный режим должен обеспечивать чтение, поиск, просмотр документов, работу с заметками и простые обновления статусов.

\subsubsection{Экран аутентификации}

Система должна предоставлять отдельный экран входа для пользователей с отображением понятных сообщений об ошибках аутентификации.

Система должна поддерживать обработку входа через внешний провайдер идентификации и отдельный режим входа для локальной разработки.

\subsubsection{Навигация по рабочим пространствам и проектам}

После входа пользователь должен получать доступ к выбору рабочего пространства, переходу к проектам и переключению между проектами в пределах доступного рабочего пространства.

\subsubsection{Главная рабочая область}

После входа пользователь должен получать доступ к основным разделам рабочего пространства: задачам, документации, файлам, заметкам, поиску и журналу активности.

\subsubsection{Интерфейс работы с задачами}

Интерфейс задач должен поддерживать фильтрацию, сортировку, поиск, изменение статуса, просмотр зависимостей, переход к связанным сущностям и отображение иерархии задач.

Интерфейс задач должен обеспечивать переход между различными представлениями задач, формируемыми на основе единой модели данных.

\subsubsection{Интерфейс базы знаний, заметок и файловых ресурсов}

Интерфейс документации должен поддерживать иерархию страниц, просмотр содержимого в формате Markdown, работу с вложенными файлами, изображениями и переходы по связанным материалам.

Интерфейс заметок и файловых ресурсов должен обеспечивать просмотр, поиск и переход к связанным задачам и документам.

\subsubsection{Интерфейс администрирования рабочего пространства}

Интерфейс администрирования должен обеспечивать управление участниками, ролями, интеграциями, агентами, учётными данными агентов и журналом аудита.

Система должна обеспечивать визуальное и логическое отделение пользовательской рабочей области от административного раздела рабочего пространства.

\subsubsection{Глобальный операторский раздел}

Система должна поддерживать отдельный раздел управления инсталляцией, не зависящий от ролей внутри конкретного рабочего пространства.

\subsubsection{Навигация и уведомления}

Система должна использовать единые навигационные элементы, хлебные крошки, уведомления об успешных операциях и понятные сообщения об ошибках.

\textheadingsep
\subsection{Требования к реализации}

\headingtextsep

Система должна быть реализована по клиент-серверной архитектуре.

Серверная часть должна быть реализована на языке Rust.

Серверная архитектура должна допускать модульную организацию приложения без обязательного разделения на отдельные микросервисы на первой версии.

Пользовательский веб-интерфейс должен быть реализован с использованием React и TypeScript.

Основное хранение структурированных данных должно выполняться в PostgreSQL.

Для полнотекстового и семантического поиска должна использоваться гибридная модель поиска на базе стандартных возможностей PostgreSQL и расширения pgvector.

Хранение файловых ресурсов должно выполняться через отдельную абстракцию хранилища, поддерживающую локальное файловое хранилище и S3-совместимое объектное хранилище.

Система должна поддерживать программный интерфейс для внешних инструментов и отдельный MCP-совместимый интерфейс для взаимодействия с агентами.

Развертывание системы должно поддерживаться в контейнеризированной среде.

\textheadingsep
\subsection{Требования к надежности}

\headingtextsep

Система должна обеспечивать сохранность пользовательских данных при штатном завершении работы и при перезапуске сервисов.

Система должна обеспечивать журналирование ошибок, предупреждений и событий аудита, достаточное для диагностики отказов.

Система должна поддерживать резервное копирование базы данных и файловых ресурсов.

При временной недоступности внешних интеграций система должна сохранять работоспособность основных внутренних функций.

После перезапуска сервисов система должна восстанавливать доступ к ранее сохранённым данным, настройкам и связям между сущностями.

\textheadingsep
\subsection{Требования к тестированию}

\headingtextsep

\subsubsection{Модульное тестирование}

Для серверной и клиентской частей должны выполняться модульные тесты ключевой бизнес-логики.

\subsubsection{Интеграционное тестирование}

Должны выполняться интеграционные тесты сценариев работы API, аутентификации, поиска, хранения данных и управления агентными учётными данными.

\subsubsection{Проверка пользовательских сценариев}

Должны выполняться сценарии проверки основных пользовательских потоков веб-интерфейса, включая вход, навигацию по проекту, работу с задачами, документами и административными разделами.

\subsubsection{Проверка миграций и восстановления}

Должны выполняться проверки применения миграций схемы данных, резервного копирования и восстановления базы данных и файловых ресурсов.

\textheadingsep
\subsection{Требования к установке}

\headingtextsep

Система должна поставляться в виде набора конфигурационных файлов и инструкций, достаточных для локального и серверного развертывания.

Установка должна поддерживать запуск в локальной среде разработки и в контейнеризированной среде.

Установка должна обеспечивать настройку подключения к PostgreSQL, выбор файлового хранилища, конфигурацию внешней аутентификации и параметров режима локальной разработки.

В состав поставки должны входить инструкции по запуску миграций, инициализации системы и настройке первого владельца рабочего пространства.

\textheadingsep
\subsection{Требования к сопровождению}

\headingtextsep

Система должна сопровождаться эксплуатационной документацией, журналами изменений и инструкциями по обновлению.

Должны выполняться процедуры миграции схемы данных, ротации учётных данных агентов, резервного копирования и восстановления системы после сбоя.

Сопровождение должно включать контроль работоспособности интеграций, журналов аудита и механизмов поиска.

\textheadingsep
\subsection{Требования к безопасности}

\headingtextsep

Система должна обеспечивать аутентификацию пользователей через внешний провайдер идентификации; в рамках данной разработки должен быть поддержан вход через GitHub OAuth при сохранении провайдер-независимой внутренней модели идентичности.

Локальный режим аутентификации для разработки не должен использоваться в рабочей конфигурации системы.

Учётные данные агентов должны храниться в виде защищённых криптографических хешей и поддерживать отзыв, ротацию, ограничение областей доступа и, при необходимости, ограничение проектом.

Веб-интерфейс должен использовать защищённые сессии, защиту от CSRF-атак и аудит критичных административных операций.

Секреты внешних интеграций должны храниться в зашифрованном виде.

Операции входа, выпуска ключей, ротации, отзыва и использования агентных учётных данных должны журналироваться и защищаться механизмами ограничения частоты запросов.

Доступ человека через веб-интерфейс и доступ агента через программный интерфейс должны рассматриваться как разные модели доступа и безопасности.

\textheadingsep
\subsection{Требования к документации}

\headingtextsep

По проекту должны быть разработаны и актуализироваться следующие документы:

\begin{itemize}
    \item техническое задание;
    \item спецификация системы;
    \item описание архитектуры системы, включая базовые архитектурные представления;
    \item материалы по прототипированию пользовательского интерфейса;
    \item сценарии работы пользователей и программных агентов;
    \item инструкция по развёртыванию и администрированию;
\end{itemize}

Документация должна вестись на русском языке и храниться в формате, пригодном для дальнейшего сопровождения и обновления.